

% ============================================================
% PDF — HOE WERK JE MET DEZE MODULE?
% ============================================================

\FVDsection{Hoe werk je met deze module?}

Deze module is in de eerste plaats een \textbf{theoriecursus}.
Begrippen, eigenschappen, redeneringen en technieken worden stap voor stap
opgebouwd en geïllustreerd met voorbeelden.

De nadruk ligt op \textbf{begrijpen}.
Het is niet voldoende dat je weet hoe je een berekening uitvoert.
Je moet ook begrijpen wat een begrip of resultaat betekent,
waarom een methode werkt en wanneer je ze kunt gebruiken.

Tussen de theorie staan op gerichte plaatsen \textbf{oefeningen}.
Die laten je meteen controleren of je een nieuw begrip of een nieuwe techniek
voldoende beheerst om verder te kunnen.
De uitgebreide inoefening en verdieping vind je in de afzonderlijke
oefeningenbundel.


% ============================================================
% BETEKENIS VAN DE LABELS — COMPACT 2 x 2
% ============================================================

\FVDsubsection{De oefeningenlabels}


% ============================================================
% RIJ 1 — ESSENTIEEL + BASIS
% ============================================================

\noindent
\begin{minipage}[t]{0.48\linewidth}

  \begin{minipage}[t]{14mm}
    \centering
    \ESS
  \end{minipage}%
  \hspace{2mm}%
  \begin{minipage}[t]{\dimexpr\linewidth-18mm\relax}

    \textbf{Essentieel}

    \vspace{1mm}

    \small
    Minimumtraject dat je zeker moet beheersen.

  \end{minipage}

\end{minipage}%
\hfill
\begin{minipage}[t]{0.48\linewidth}

  \begin{minipage}[t]{14mm}
    \centering
    \oefB
  \end{minipage}%
  \hspace{2mm}%
  \begin{minipage}[t]{\dimexpr\linewidth-18mm\relax}

    \textbf{Basis}

    \vspace{1mm}

    \small
    Noodzakelijke basisvaardigheden en standaardtechnieken.

  \end{minipage}

\end{minipage}


\vspace{6mm}


% ============================================================
% RIJ 2 — VERDIEPING + GEVORDERD
% ============================================================

\noindent
\begin{minipage}[t]{0.48\linewidth}

  \begin{minipage}[t]{14mm}
    \centering
    \oefU
  \end{minipage}%
  \hspace{2mm}%
  \begin{minipage}[t]{\dimexpr\linewidth-18mm\relax}

    \textbf{Verdieping}

    \vspace{1mm}

    \small
    Meer inzicht, verbanden en transfer naar nieuwe situaties.

  \end{minipage}

\end{minipage}%
\hfill
\begin{minipage}[t]{0.48\linewidth}

  \begin{minipage}[t]{14mm}
    \centering
    \oefG
  \end{minipage}%
  \hspace{2mm}%
  \begin{minipage}[t]{\dimexpr\linewidth-18mm\relax}

    \textbf{Gevorderd}

    \vspace{1mm}

    \small
    Hogere complexiteit, modelleren, redeneren of bewijzen.

  \end{minipage}

\end{minipage}


% ============================================================
% VERDUIDELIJKING — ICONEN IN DE ZIN
% ============================================================

\vspace{5mm}

\noindent
\small
\textbf{Let op:}\quad
\ESSinline\ staat los van
\oefBinline, \oefUinline\ en \oefGinline:
ook een verdiepende of gevorderde oefening kan essentieel zijn.

\par
\medskip


% ============================================================
% ACTIEF STUDEREN
% ============================================================

\FVDsubsection{Werk actief met de theorie}

Lees voorbeelden niet alleen door, maar probeer tussenstappen zelf te
verklaren. Vraag jezelf regelmatig af waarom een bepaalde werkwijze geldig is
en wat het verkregen resultaat betekent.

Maak de oefeningen met de diamant wanneer je ze tegenkomt.
Controleer daarbij niet alleen of je het juiste antwoord vindt.
Je moet je werkwijze ook kunnen uitleggen, je keuzes kunnen verantwoorden
en, waar nodig, je resultaat kunnen interpreteren.

Bij toepassingen betekent dit dat je ook de stap moet kunnen maken van een
concrete situatie naar een wiskundig model en van het wiskundige resultaat
terug naar de oorspronkelijke context.



